\chapter{Hamiltonian Monte Carlo}\label{ch:hmc}

The random walk Metropolis--Hastings algorithm explored in
\cref{ch:bayesian} generates proposals by perturbing the current state
with isotropic Gaussian noise. This strategy ignores the geometry of the
posterior: proposals are equally likely in every direction, regardless of
whether the posterior is steep or flat. Now that we have the ability to
differentiate through the forward map, we can exploit gradient information
to guide the sampler along the posterior surface. \emph{Hamiltonian Monte
Carlo} (HMC)~\cite{betancourt2018hamiltonian,neal2011mcmc} does precisely
this by borrowing the formalism of Hamiltonian dynamics to construct
proposals that are both distant from the current state and likely to be
accepted.

The physical intuition is as follows. We imagine the current position of
the Markov chain as a particle moving in the parameter space. At each
step, the particle is given a random kick (a momentum draw) and then
allowed to evolve under a Hamiltonian that couples its position to the
posterior. Because the dynamics conserve energy, the particle naturally
spends more time in regions of high posterior density and traces out
long-range, low-energy trajectories that would take a random walk many
steps to traverse.

\section{The Hamiltonian}\label{sec:hamiltonian}

Let $q \in \mathbb{R}^d$ denote the position (i.e.\ the parameter
vector~$\phi$; we use $q$ to follow the Hamiltonian mechanics convention) and
introduce an auxiliary momentum variable $p \in \mathbb{R}^d$. We define
the joint density of position and momentum through the canonical
distribution
\begin{equation}
  P(q, p) \propto \exp\!\big({-H(q, p)}\big),
  \label{eq:canonical-distribution}
\end{equation}
where $H(q, p)$ is the Hamiltonian. We decompose $H$ into a potential
energy~$U(q)$ and a kinetic energy~$K(p)$:
\begin{equation}
  H(q, p) = U(q) + K(p).
  \label{eq:hamiltonian-decomposition}
\end{equation}
This separation ensures that the joint density factorises as
$P(q, p) \propto \exp(-U(q))\,\exp(-K(p))$, so the position and
momentum are independent: marginalising over~$p$ recovers the target
distribution over~$q$ alone.

\subsection{Potential Energy}

We define the potential energy as the negative log-posterior:
\begin{equation}
  U(q) = -\log p(q \mid \mathbf{y}) = -\ell(q) + \text{const},
  \label{eq:potential-energy}
\end{equation}
where the constant absorbs the normalising factor. Under this
identification, regions of high posterior density correspond to low
potential energy. The physical picture is then that of a landscape in
which the particle rolls into valleys at the posterior modes and must
climb uphill to escape --- precisely the behaviour we want from a
sampler.

\subsection{Kinetic Energy}

The standard choice of kinetic energy is the quadratic form
\begin{equation}
  K(p) = \tfrac{1}{2}\, p^T M^{-1} p,
  \label{eq:kinetic-energy}
\end{equation}
where $M \in \mathbb{R}^{d \times d}$ is a symmetric positive-definite
\emph{mass matrix}. Under~\cref{eq:canonical-distribution} the momentum
is then distributed as $p \sim \mathcal{N}(0, M)$. The mass matrix
controls the inertia of the particle in each direction: a well-chosen~$M$
that approximates the posterior covariance can precondition the dynamics
so that the particle explores all directions at a similar rate.

\section{Algorithm}\label{sec:hmc-algorithm}

Hamilton's equations of motion for the
system~\cref{eq:hamiltonian-decomposition}
are
\begin{equation}
  \frac{dq}{dt} = \frac{\partial H}{\partial p} = M^{-1}p,
  \qquad
  \frac{dp}{dt} = -\frac{\partial H}{\partial q} = -\nabla_q U(q).
  \label{eq:hamiltons-equations}
\end{equation}
The gradient $\nabla_q U(q) = -\nabla_q \ell(q)$ is the
derivative of the log-likelihood with respect to the parameters, computed
via the adjoint method described in \cref{ch:adjoint}.

\subsection{Leapfrog Integration}

In practice Hamilton's equations are integrated numerically using the
\emph{leapfrog} integrator, which is symplectic and
time-reversible. Given a step size~$\varepsilon$ and $L$~integration
steps, one leapfrog trajectory from $(q_0, p_0)$ proceeds as:
\begin{equation}
  \begin{aligned}
    p_{i+1/2} &= p_i - \tfrac{\varepsilon}{2}\,\nabla_q U(q_i), \\
    q_{i+1}   &= q_i + \varepsilon\, M^{-1} p_{i+1/2}, \\
    p_{i+1}   &= p_{i+1/2} - \tfrac{\varepsilon}{2}\,\nabla_q U(q_{i+1}),
  \end{aligned}
  \label{eq:leapfrog}
\end{equation}
for $i = 0, 1, \ldots, L-1$. Each step requires one evaluation of
$\nabla_q U$, so a trajectory of $L$~steps costs $L$~gradient
evaluations, each of which involves a full forward and adjoint solve of
the model.

\subsection{Metropolis Correction}

Let $(q^*, p^*)$ denote the endpoint of the leapfrog trajectory starting
from $(q, p)$. Because the leapfrog integrator introduces a discretisation
error, the Hamiltonian is not exactly conserved along the trajectory. To
ensure the chain targets the correct stationary distribution, a
Metropolis acceptance step is applied. The sign of the momentum at the
endpoint is first negated, $p^* \to -p^*$, to make the proposal a
time-reversal of the forward trajectory; this is necessary to satisfy the
detailed balance condition. The acceptance probability is then
\begin{equation}
  r = \min\!\big\{1,\; \exp\!\big(H(q, p) - H(q^*, p^*)\big)\big\}.
  \label{eq:hmc-acceptance}
\end{equation}
Because the leapfrog integrator is symplectic, the energy error is
typically small and the acceptance rate is high even for proposals that
are far from the current state.

Note that after the accept/reject decision is made the momentum is
discarded; the position~$q$ (accepted or not) is retained as the next
sample of the parameter vector. At the start of the following iteration a
fresh momentum is drawn from $\mathcal{N}(0, M)$, which randomises the
energy level and ensures the particle explores the full posterior rather
than tracing a single energy contour throughout.

The full HMC procedure is summarised in \cref{alg:hmc}.

\begin{algorithm}
  \caption{Hamiltonian Monte Carlo}\label{alg:hmc}
  \begin{algorithmic}[1]
    \Require Initial state $q^{(0)}$, mass matrix~$M$, step
    size~$\varepsilon$, number of leapfrog steps~$L$, number of
    iterations~$T$
    \For{$t = 0, 1, \ldots, T-1$}
    \State Draw $p \sim \mathcal{N}(0, M)$
    \State Set $(q_0, p_0) \gets (q^{(t)}, p)$
    \For{$i = 0, 1, \ldots, L-1$}
    \State Perform one leapfrog step~\cref{eq:leapfrog}:
    $(q_i, p_i) \to (q_{i+1}, p_{i+1})$
    \EndFor
    \State Set $(q^*, p^*) \gets (q_L, -p_L)$
    \Comment{Negate momentum}
    \State Compute $\Delta H = H(q^{(t)}, p) - H(q^*, p^*)$
    \State Draw $u \sim \mathrm{Uniform}(0,1)$
    \If{$u < \min\{1,\, \exp(\Delta H)\}$}
    \State $q^{(t+1)} \gets q^*$
    \Else
    \State $q^{(t+1)} \gets q^{(t)}$
    \EndIf
    \EndFor
    \Ensure Chain $\{q^{(0)}, q^{(1)}, \ldots, q^{(T)}\}$
  \end{algorithmic}
\end{algorithm}

\section{Tuning the Hyperparameters}\label{sec:hmc-tuning}

The efficiency of HMC depends on three hyperparameters: the step
size~$\varepsilon$, the number of leapfrog steps~$L$, and the mass
matrix~$M$. Poor choices lead to either high rejection rates (if
$\varepsilon$ is too large), slow exploration (if $L$ is too small), or
wasted computation (if $L$ is too large and the trajectory doubles back on
itself).

The most widely used adaptive scheme is the No-U-Turn Sampler
(NUTS)~\cite{hoffman2014no}, which eliminates the need to specify~$L$ by
dynamically extending the trajectory until a U-turn criterion is met.
While effective in general, NUTS can require many leapfrog steps per
sample --- each of which, in our setting, demands a full forward and
adjoint solve of the electrochemical model.

In this work we instead use the \emph{ChEES-HMC} criterion of
Hoffman et al.~\cite{hoffman2021adaptive}, which adapts $\varepsilon$,
$L$, and~$M$ jointly during a warm-up phase by maximising the
\emph{Change in the Estimator of the Expected Square} (ChEES) of the
trajectory length. We found that ChEES-HMC consistently selected shorter
trajectory lengths than NUTS for this problem whilst maintaining
comparable effective sample sizes. Since each leapfrog step requires a
gradient evaluation --- and each gradient evaluation requires a full
forward and adjoint solve --- reducing the number of integration steps
per sample directly reduces the computational cost. This makes ChEES-HMC
particularly well suited to scientific applications in which the gradient
of the forward model is expensive to evaluate.

% TODO: I have not introduced effective sample size yet

During warm-up, the step size and trajectory length are adapted alongside
the mass matrix, and the sampler simultaneously performs gradient descent
toward the posterior mode, serving the dual purpose of burn-in and
hyperparameter tuning. Once the warm-up phase concludes, the
hyperparameters are fixed and all subsequent samples are retained as draws
from the posterior.

\begin{tabular}{lcc}
  Parameter & HMC $\hat{R}$ & RWMH $\hat{R}$ \\
  $\alpha$ & 1.0133 & 1.0106 \\
  $K_0$ & 1.0013 & 1.0010 \\
  $\theta_f$ & 1.0015 & 1.0015 \\
\end{tabular}
